\documentclass[a4paper,12pt]{article}
\usepackage[T2A]{fontenc}
\usepackage[utf8]{inputenc}
\usepackage[russian]{babel}
\usepackage{amsmath,amssymb}
\usepackage{helvet}
\renewcommand{\familydefault}{\sfdefault}
\usepackage{geometry}
\geometry{left=22mm,right=22mm,top=22mm,bottom=22mm}
\usepackage{microtype}
\usepackage{tabularx,booktabs,array}
\usepackage{enumitem}
\usepackage{fancyhdr}
\usepackage{parskip}
\usepackage{xcolor}
\usepackage{tikz}
\usepackage{hyperref}
\hypersetup{hidelinks}
\definecolor{accent}{HTML}{008080}
\definecolor{deepaccent}{HTML}{005D5D}
\definecolor{textgray}{HTML}{353C3C}
\definecolor{pale}{HTML}{E8F3F1}
\color{textgray}
\setlength{\parindent}{0pt}
\setlength{\parskip}{5pt}
\setlength{\headheight}{15pt}
\pagestyle{fancy}
\fancyhf{}
\fancyhead[L]{\footnotesize МАТЕМАТИКА\quad\textcolor{accent}{|}\quad Ярослав Верещагин}
\fancyhead[R]{\footnotesize 12.11.2026}
\fancyfoot[L]{\footnotesize\textcolor{accent}{Лёвин А.~А.}}
\fancyfoot[R]{\footnotesize\thepage}
\renewcommand{\headrulewidth}{0pt}
\renewcommand{\footrulewidth}{0pt}
\newcommand{\Cover}[3]{%
  \thispagestyle{empty}
  \begin{center}
  \vspace*{20mm}
  {\large\color{deepaccent} МАТЕМАТИКА\par}
  \vspace{21mm}
  {\LARGE #1\par}
  \vspace{11mm}
  {\Large\color{accent} #2\par}
  \vspace{20mm}
  {\large Подготовка к ОГЭ по математике\par}
  \vspace{40mm}
  {\large Ярослав Верещагин\par}
  \vspace{15mm}
  {\large 12 ноября 2026 г.\par}
  \vfill
  {\small Лёвин Артём Александрович\par}
  {\small эксклюзивно для Ярослава Верещагина\par}
  \vspace{15mm}
  \begin{tikzpicture}
     \draw[color=accent, line width=.75pt] (0,0) .. controls (1.3,.55) and (2.8,-.42) .. (4.0,0)
      .. controls (5.0,.47) and (6.1,-.35) .. (7.1,0)
      .. controls (8.1,.41) and (9.1,-.48) .. (10.5,0);
  \end{tikzpicture}
  \end{center}
  \clearpage
}
\newcommand{\Section}[1]{\vspace{2mm}{\large\bfseries\color{deepaccent}#1}\par\vspace{2mm}}
\newcommand{\Closing}{\vspace{8mm}\begin{center}{\footnotesize\color{deepaccent}Лёвин Артём Александрович\par эксклюзивно для Ярослава Верещагина}\end{center}}
\begin{document}

\Cover{Проверка домашней работы}{Задания чек-листа\\[4pt]\large Движение по реке}{}
\Section{Общий результат}
Сверены решения на трёх фотографиях с десятью заданиями чек-листа «Движение по реке» от 9 ноября 2026 г.

\begin{tabularx}{\linewidth}{@{}p{66mm}X@{}}
\toprule
\textbf{Показатель} & \textbf{Результат} \\
\midrule
Представлено решений & 6 из 10 \\
Подтверждённые вычислительные ошибки & Не обнаружены \\
Недостающие письменные пояснения / проверки & № 2, 3, 4, 5 \\
Не представлены & № 7, 8, 9, 10 \\
\bottomrule
\end{tabularx}

В представленных задачах \textbf{конечные ответы правильные}. Следующие замечания относятся к отдельным требованиям условий, которые полезно дописать. В нескольких заданиях требуется объяснение или обратная подстановка: одного численного ответа для полного оформления недостаточно.

\Section{Неполные элементы решений и закрепление}
\textbf{№ 2. Плот: требуется пояснить скорость.}

В условии отдельно просят объяснить, почему для свободно плывущего плота не вводится собственная скорость. Плот движется вместе с водой без работающего двигателя. Поэтому его скорость относительно берега совпадает со скоростью течения. Именно эту скорость нужно использовать в формуле \(t=S/v\).

\textit{Закрепление.} Плот отпустили на реке со скоростью течения \(2{,}5\) км/ч. Какое время понадобится плоту для прохождения 20 км? Запишите, какая скорость используется. \textbf{Ответ:} 8 ч; скорость течения \(2{,}5\) км/ч.

\medskip
\textbf{№ 3. Две скорости: требуется явная проверка.}

После нахождения собственной скорости и скорости течения условие просит проверить обе исходные скорости. Проверка должна содержать два действия: \(15+3=18\) км/ч по течению и \(15-3=12\) км/ч против течения. Это подтверждает, что оба наблюдения выполнены одновременно.

\textit{Закрепление.} Лодка идёт по течению со скоростью 22 км/ч, против течения --- 14 км/ч. Определите обе неизвестные скорости и проверьте исходные значения. \textbf{Ответ:} 18 км/ч и 4 км/ч; \(18+4=22\), \(18-4=14\).

\medskip
\textbf{№ 4. Разность времени: требуется объяснение знака.}

Для одной и той же длины пути движение против течения дольше, поскольку \(x-2<x+2\) при \(x>2\). При положительных скоростях меньшая скорость даёт большее время: \(\dfrac{30}{x-2}>\dfrac{30}{x+2}\). Поэтому из большего времени вычитается меньшее. В работе уравнение составлено верно; это объяснение следует записать словами.

\textit{Закрепление.} Лодка со своей скоростью 10 км/ч идёт 32 км против течения и столько же по течению. Течение --- 2 км/ч. Найдите времена и укажите, какой переход дольше. \textbf{Ответ:} 4 ч и 2 ч 40 мин; дольше против течения.

\clearpage
\textbf{№ 5. Скорость течения: требуется проверка по времени.}

Найденную скорость течения необходимо подставить в исходное условие. Вниз по реке скорость катера составляет \(10+2=12\) км/ч, вверх --- \(10-2=8\) км/ч. Времена: \(48:12=4\) ч и \(48:8=6\) ч; разность \(6-4=2\) ч, как в условии. Неподходящий отрицательный корень исключён правильно.

\textit{Закрепление.} Собственная скорость судна равна 12 км/ч, течение --- 3 км/ч. Судно проходит 60 км по течению и 60 км против течения. Найдите оба времени и их разность. \textbf{Ответ:} 4 ч, 6 ч 40 мин; разность 2 ч 40 мин.

\Section{Оставшаяся часть чек-листа}
\begin{tabularx}{\linewidth}{@{}>{\bfseries}p{11mm} >{\raggedright\arraybackslash}p{37mm} X@{}}
\toprule
№ & Тема & Первый шаг \\
\midrule
7 & Рейс по расписанию & Общее время --- 10 ч; после часовой стоянки на движение остаётся 9 ч. \\
8 & Плот и догоняющая лодка & Плот движется 9 ч, моторная лодка --- 7 ч. \\
9 & Неизвестная скорость реки & После исключения часовой стоянки движение занимает 10 ч. Скорости: \(15+u\) и \(15-u\). \\
10 & Задержка и остановка & Плот движется 12 ч, теплоход отправился на 2 ч позже и стоял 1 ч. На движение теплохода остаётся 9 ч. \\
\bottomrule
\end{tabularx}

Для каждого оставшегося номера составьте таблицу «скорость --- время --- расстояние», затем одно уравнение по сумме или разности времён. Проверьте положительность скорости против течения и подставьте ответ в исходное условие.

\Section{Итог}
Численные результаты шести представленных заданий подтверждены. Для полного оформления нужно дописать указанные объяснения и проверки; для завершения чек-листа --- решить задачи № 7--10.
\Closing
\end{document}
